\section{DQN 方法与实验}
\label{sec:dqn}
本节按算法原理、模型实现、训练设置与开发分析的顺序介绍 DQN。除明确引用的统一测试外，本节分数均属于训练开发评估或开发复测。

\subsection{算法原理}
\subsubsection{动作价值与 Q-learning}
动作价值函数描述在状态 $s$ 下先执行动作 $a$，随后按策略 $\pi$ 行动时的预期回报：
\begin{equation}
 Q^\pi(s,a)=\mathbb{E}_{\pi}[G_t\mid s_t=s,a_t=a].
\end{equation}
在同一画面下，向左或向右移动可能带来不同的后续接球机会。价值函数将这种长期差异转化为可比较的数值。若最优动作价值函数 $Q^*$ 已知，就可以采用 $\arg\max_a Q^*(s,a)$ 选择动作。

最优动作价值满足 Bellman 最优关系，即当前动作的价值由即时奖励和下一状态的最优未来价值共同决定：
\begin{equation}
 Q^*(s,a)=\mathbb{E}\!\left[
 r+\gamma(1-d)\max_{a'}Q^*(s',a')\mid s,a
 \right],
\end{equation}
其中 $s'$ 是下一状态，$d$ 表示真正终止。表格型 Q-learning 利用一次环境转移更新对应表项：
\begin{equation}
 Q(s_t,a_t)\leftarrow Q(s_t,a_t)+\alpha
 \left[r_t+\gamma(1-d_t)\max_{a'}Q(s_{t+1},a')-Q(s_t,a_t)\right].
\end{equation}
这里 $\alpha$ 为学习率，方括号中的差值是时间差分误差。更新通过下一状态的估计值修正当前价值，因此不必等一局游戏全部结束后才开始学习。

图像状态的组合数量远大于表格能够存储的范围。DQN 用参数为 $\theta$ 的神经网络 $Q_\theta(s,a)$ 近似动作价值，并通过共享视觉特征将已见画面的经验迁移到相似画面。经典 DQN 将卷积表示与价值学习结合，建立了从像素输入到离散动作决策的学习框架。本项目沿用该框架，网络输出四个动作的 Q 值，而不是直接输出动作概率。

\subsubsection{经验回放与目标网络}
神经网络训练通常希望样本具有较低的相关性，而游戏中相邻画面非常相似。若严格按采样顺序立即更新，参数容易受近期轨迹支配。经验回放池 $\mathcal D$ 保存转移
\begin{equation}
 e_t=(s_t,a_t,r_t,s_{t+1},d_t),
\end{equation}
训练时从池中随机抽取小批量样本。这样既减弱相邻转移之间的相关性，也允许同一条有效经验被多次利用。低奖励和失球经验同样保留在池中，因为区分成功与失败对动作评价都具有意义。

DQN 同时维护在线网络 $Q_\theta$ 与目标网络 $Q_{\theta^-}$。在线网络参与梯度更新并用于决策；目标网络在一段时间内保持不变，用于构造学习目标，随后从在线网络复制参数。一步 TD 目标和误差写为
\begin{align}
 y_t &= r_t+\gamma(1-d_t)\max_{a'}Q_{\theta^-}(s_{t+1},a'),\\
 \delta_t &= y_t-Q_\theta(s_t,a_t).
\end{align}
目标网络将监督目标的变化与每次在线更新适度分离，缓解自举估计与参数更新相互追逐的问题。计算 $y_t$ 时停止梯度传播，梯度仅作用于在线网络。

本项目使用 Stable-Baselines3 的 DQN 实现，采用 Huber 损失：
\begin{equation}
 \ell(\delta)=
 \begin{cases}
 \frac{1}{2}\delta^2,& |\delta|\leq 1,\\
 |\delta|-\frac{1}{2},& |\delta|>1,
 \end{cases}
 \qquad
 \mathcal L(\theta)=\frac{1}{B}\sum_{i=1}^{B}\ell(\delta_i).
\end{equation}
其中 $B$ 为批大小。小误差区域使用平方形式，大误差区域使用线性形式，从而减弱个别较大 TD 误差对更新的影响。优化器采用 Adam。损失减小表示网络更接近当前样本上的 TD 目标，并不直接意味着下一次评估的游戏得分必然提高。

\subsubsection{探索机制与多步回报}
训练采用 $\varepsilon$-greedy 策略：
\begin{equation}
 a_t=
 \begin{cases}
 \text{从动作空间均匀采样},&\text{概率为 }\varepsilon_t,\\
 \arg\max_a Q_\theta(s_t,a),&\text{概率为 }1-\varepsilon_t.
 \end{cases}
\end{equation}
早期较大的 $\varepsilon_t$ 帮助收集不同状态与动作的经验，后期降低探索率以更多利用已学策略。开发评估采用贪心动作，因此训练阶段含随机探索的回合分数与评估分数应分别解释。

对于延迟奖励，一步目标需要通过多次自举将价值向前传播。项目还尝试了三步回报。中间没有真正终止时，其目标为
\begin{equation}
 y_t^{(3)}=r_t+\gamma r_{t+1}+\gamma^2r_{t+2}
 +\gamma^3\max_{a'}Q_{\theta^-}(s_{t+3},a').
\end{equation}
若提前终止，则仅累积终止前的奖励，并去除终止后的自举项。多步回报可以更直接地利用后续实际奖励，但也引入更长轨迹中的随机性和行为策略影响，其效果需要结合实验判断。本项目主要结果采用一步 DQN，三步设置作为后续恢复实验。

\subsection{模型实现与训练管理}
\subsubsection{卷积价值网络}
模型采用 \texttt{CnnPolicy}，使用 NatureCNN 风格的卷积特征提取器。输入的四帧图像沿通道维拼接，卷积同时提取空间结构与跨帧差异。三层卷积之后将特征展平，经过 512 维全连接层，再由动作价值输出层产生四个 Q 值。各层使用的尺寸列于表\ref{tab:network}。

\begin{table}[!htbp]
 \centering\small
 \caption{DQN 卷积价值网络结构，省略批维度}\label{tab:network}
 \begin{tabularx}{\linewidth}{l Y l}
 \toprule
 层 & 运算与参数 & 输出尺寸\\
 \midrule
 输入 & 四帧灰度图像；像素除以 255 & $4\times84\times84$\\
 卷积 1 & 32 个 $8\times8$ 卷积核，步长 4，ReLU & $32\times20\times20$\\
 卷积 2 & 64 个 $4\times4$ 卷积核，步长 2，ReLU & $64\times9\times9$\\
 卷积 3 & 64 个 $3\times3$ 卷积核，步长 1，ReLU & $64\times7\times7$\\
 展平 & 按通道和空间位置展开 & 3136\\
 全连接 & $3136\rightarrow512$，ReLU & 512\\
 输出层 & $512\rightarrow4$，线性输出 & 4\\
 \bottomrule
 \end{tabularx}
\end{table}

输出层不使用 softmax，因为 Q 值表示预期回报，可以是任意实数。动作选择依据输出值的相对大小进行。在线网络与目标网络具有相同结构，目标网络由在线网络周期性复制得到。当前实现使用普通 DQN 的最大目标值计算，没有加入 Double DQN、Dueling 结构或优先经验回放。

网络不显式预测小球落点，也不包含人工编写的追球控制规则。像素特征是否有利于接球，需要通过环境交互产生的奖励逐步学习。由此，模型实现将视觉表示与控制目标连接起来，而游戏的物理状态仅用于环境内部更新和调试。

\subsubsection{训练循环与模型更新}
DQN 训练流程如下，其中环境步均指经过动作重复包装后的决策步：
\begin{enumerate}
 \item 初始化在线网络，复制其参数得到目标网络，建立回放池并重置环境。
 \item 根据当前探索率选择动作，执行环境步，获得新观测、累计奖励和结束信息。
 \item 将转移存入回放池；回合结束后重置环境，继续收集完整游戏经验。
 \item 超过预热步数后，每收集 4 个决策步执行 1 次梯度更新，每次随机采样 32 条转移。
 \item 使用目标网络构造 TD 目标，以 Huber 损失更新在线网络；每 10000 个决策步同步目标网络。
 \item 定期运行独立开发评估，保存检查点和表现最好的模型，达到预算或早停条件后结束。
\end{enumerate}

这里动作重复 4 帧与每 4 个决策步训练一次是两个不同层次的设置。前者控制智能体多久重新作出一次决策，后者控制采样与参数更新的比例。在未提前结束的情况下，一次梯度更新通常对应最近约 16 个物理帧的采样。报告中的 50 万步等训练预算均采用决策步计数。

\subsubsection{续训与回放状态保存}
模型文件保存网络参数、优化器状态及相关配置，但不自动包含经验回放池。项目因此把回放池作为独立产物保存，并检查其环境定义、容量与回报跨度是否与待恢复模型兼容。压缩用于降低磁盘占用，训练时解压后的回放仍需要相应内存。

在缺少配套回放池的续训中，模型从已有权重开始建立新的回放池。项目实现了策略预热：前 20000 个决策步仍暂缓梯度更新，但动作由恢复后的 $\varepsilon$-greedy 策略产生。这样，预热数据与当前策略的访问状态更接近，而不是重新由完全随机动作填充。这种操作属于基于已有模型的新阶段训练，不等同于完整恢复原训练过程。

DQN 训练入口为 \texttt{train.py}，策略预热位于 \texttt{warm\_start\_dqn.py}，回放兼容性与压缩恢复由 \texttt{replay\_io.py} 实现。

\subsection{训练设置与日志诊断}
\subsubsection{基础训练设置}
DQN 基础实验在固定随机性设置下，从随机初始化开始训练 500000 个决策步。初始探索率为 1.0，在前 60\% 的训练预算内线性下降至 0.05，之后保持不变。学习率固定为 $10^{-4}$，回放容量为 50000，预热步数为 20000。主要参数与续训参数的对照见表\ref{tab:hyper}。

\begin{table}[!htbp]
 \centering\small
 \caption{从头训练与当前低学习率续训配置}\label{tab:hyper}
 \begin{tabularx}{\linewidth}{Y c c}
 \toprule
 参数 & 从头训练 & 低学习率续训\\
 \midrule
 阶段预算（决策步） & 500000 & 200000\\
 学习率 & $1\times10^{-4}$ & $3\times10^{-5}\rightarrow1\times10^{-5}$\\
 折扣因子 $\gamma$ & 0.99 & 0.99\\
 回放池容量 & 50000 & 50000\\
 新回放池预热步数 & 20000 & 20000\\
 批大小 & 32 & 32\\
 训练频率 / 每次梯度步数 & 4 / 1 & 4 / 1\\
 目标网络同步间隔 & 10000 & 10000\\
 探索率变化 & $1.00\rightarrow0.05$ & $0.05\rightarrow0.02$\\
 探索率衰减占阶段比例 & 60\% & 100\%\\
 回报跨度 & 1 步 & 1 步\\
 评估间隔（决策步） & 50000 & 25000\\
 每次评估局数 & 20 & 50\\
 \bottomrule
 \end{tabularx}
\end{table}

图像转移需要较大的内存。按当前回放规模估算，分别以 \texttt{uint8} 保存当前与下一堆帧观测时，仅观测约占 2.8 GB，其他字段和运行过程还会带来额外开销。因此，回放容量同时影响可保留的历史范围与内存需求。

训练过程记录回合奖励、回合长度、探索率、损失、回放池大小和梯度更新次数等信息，并定期保存模型检查点。基础实验每 50000 步进行一次 20 局开发评估，使用评估均分选取最佳模型。

\subsubsection{低学习率续训与策略保护}
基础模型获得一定控制能力后，项目开展低学习率续训，尝试在已有策略附近继续改进。续训阶段预算为 200000 步，学习率从 $3\times10^{-5}$ 线性下降到 $10^{-5}$，探索率从 0.05 下降到 0.02，每 25000 步评估 50 局。减小学习率有助于缩小单次参数变化幅度，但能否维持已有能力仍取决于新经验分布和价值估计。

现有记录显示，低学习率续训在阶段早期取得最佳开发均分，此后继续训练并未保持这一水平。针对该现象，当前续训配置加入起点评估、策略预热和连续多次开发评估未改进时早停等保护。起点评估先把恢复的模型纳入候选，避免一次较弱的新检查点被当成该阶段的最佳模型；预热期间尚未更新参数，重复评估不会消耗早停耐心；产生新最佳结果时同步保存模型和对应回放池。

上述保护是工程整理后的当前配置。历史低学习率实验完整运行至 200000 步，其结果不能被解释为当前早停配置下重新运行的结果。配置用于说明可复现的训练入口，历史表现则按项目已有记录报告。

\subsubsection{后续恢复实验与评估协议}
在低学习率续训之后，项目又尝试了三步回报恢复和一步扩大回放池两种设置。前者希望加快延迟奖励的利用，后者希望延长历史经验的保留时间。两次实验均未超过恢复起点的开发均分，并提前停止。历史配置快照显示，两项恢复实验采用相同的扩大回放容量；它们从同一个续训最佳模型分别开始，而非彼此顺序衔接。

开发评估在独立环境中进行，使用贪心策略，且不更新参数。训练回调在固定随机性设置下生成连续游戏序列，模型选择所用数据属于开发数据。基础训练每次评估 20 局，续训与恢复实验每次评估 50 局，样本规模不同。

独立复测采用逐局重置的方式，使不同候选模型面对相同的开局集合，共运行 30 局。它与训练回调的连续游戏序列是两套不同协议，各自的均分分别报告，不混合计算。

\reportfigure{images/training_curves.png}{训练与开发评估曲线}{上排为训练回合奖励，下排为贪心开发评估均分；各列独立计步。}{四个实验阶段的实际训练日志。上排为全部回合得分及滚动均值，下排为各评估节点均分；星号标出最佳节点，虚线标出两项恢复实验的早停，浅灰背景表示续训新回放池的预热区间。}{fig:curves}

\subsubsection{探索率与损失日志诊断}
为进一步观察训练信号，图\ref{fig:diagnostics}直接读取两个主要阶段的训练日志，分别绘制探索率和 Huber 损失。日志中的缺失损失值不插值，平滑曲线采用有效日志点的滚动中位数；由于日志记录间隔不完全一致，该窗口不代表固定的决策步数。

从头训练的探索率降至设定下限后，开发均分仍出现明显回落，说明后期回退不能仅用探索率继续下降来解释。低学习率续训的探索率缓慢下降，最佳结果却出现在续训早期，也说明探索安排和策略得分之间没有简单的单调关系。

损失同样不能作为选模的替代指标。基础训练早期的损失较小，后期有所增大，但较小的拟合误差并未对应较强的游戏控制。随着模型访问状态和回放样本变化，TD 目标也会变化，因此损失的尺度和波动必须结合当前经验分布理解。这些日志支持使用独立开发得分选取模型，但尚不足以单独识别回退的具体原因。

\reportfigure{images/training_diagnostics.png}{探索率与损失诊断}{基础训练和低学习率续训的实际探索率、Huber 损失日志。}{两项主要实验的探索率与损失变化。上排为探索概率，下排为对数轴下的原始损失点及滚动中位数；红色虚线标记开发最佳节点。预热期间没有损失记录的区间保留为空白。}{fig:diagnostics}

\subsection{开发结果与行为分析}
各阶段结果及训练曲线来自训练监控与开发评估日志，保留全部实际采样点。为补齐配对分布图，本次固定已有权重，在 CPU 上对两个候选模型分别运行相同开局的 30 局贪心评估；均分复现原记录的 260.27 与 122.40。逐局分数、原始行为帧和观测样本保存在项目的实验数据目录中，这些复测仍属于开发数据。

\subsubsection{各训练阶段的结果}
表\ref{tab:results}给出当前物理版本下 DQN 各阶段的实验结果。表中最佳节点指该阶段开发评估均分最高的节点；后两个恢复实验的最佳模型均为恢复起点，而非新训练得到的模型。

\begin{table}[!htbp]
 \centering\small
 \caption{各阶段开发评估结果}\label{tab:results}
 \begin{tabularx}{\linewidth}{Y P{3.0cm} r r c}
 \toprule
 实验阶段 & 最佳模型所在阶段 & 最佳均分 & 结束均分 & 局数\\
 \midrule
 从头训练 500k & 训练中后期 & 174.00 & 133.30 & 20\\
 一步低学习率续训 200k & 续训早期 & 283.48 & 113.12 & 50\\
 三步回报恢复 & 恢复起点 & 283.48 & 187.78 & 50\\
 一步扩大回放 & 恢复起点 & 283.48 & 216.94 & 50\\
 \bottomrule
 \end{tabularx}
 \vspace{.3em}
 \begin{minipage}{\linewidth}\footnotesize
 注：$1\mathrm{k}=1000$ 个决策步；“局数”为每次评估局数。不同阶段的评估规模不同，表中数值用于呈现各阶段轨迹与模型选择结果。
 \end{minipage}
\end{table}

基础训练在中后期取得 174.00 的最佳均分，结束时下降到 133.30，减少 40.70 分。低学习率续训的最佳均分为 283.48，出现在续训早期；结束时为 113.12，比最佳节点低 170.36 分。两组结果共同说明，增加训练步数并不保证策略持续改善，最终模型与最佳模型需要分别保存。

基础阶段和续训阶段的最高均分来自不同局数的评估，并且续训同时改变了模型起点、训练预算、学习率与探索安排。因此，不能把二者的差值直接归因于降低学习率。更有信息量的比较来自同一协议下的候选模型复测。

\subsubsection{相同开发种子下的候选模型复测}
在相同开局的 30 局开发复测中，当前推荐的最佳模型取得平均分 260.27、中位数 257.00；上一轮最终模型的平均分为 122.40，如表\ref{tab:audit}所示。

\begin{table}[!htbp]
 \centering
 \caption{相同 30 个开发种子上的独立复测}\label{tab:audit}
 \begin{tabularx}{.90\linewidth}{Y r r c}
 \toprule
 候选模型 & 平均分 & 中位数 & 局数\\
 \midrule
 当前推荐最佳模型 & 260.27 & 257.00 & 30\\
 上一轮最终模型 & 122.40 & 101.50 & 30\\
 \bottomrule
 \end{tabularx}
\end{table}

两者平均分相差 137.87 分，相对增加约 112.64\%。由于这次比较使用相同种子集合，它支持在所测开局条件下选择当前推荐模型。最佳模型在 28/30 个种子上得分更高；其得分范围为 36--444 分，仍存在明显的局间波动。箱线图与逐局散点共同呈现这种差异，不能仅凭均值和中位数接近就推断策略稳定。

本次配对复测使用同一低学习率续训阶段保存的最佳检查点与结束检查点。两者均固定权重进行推理，具体模型与逐局分数可由项目中的评估记录核对。

\reportfigure{images/score_distribution.png}{相同开局下的得分分布}{30 局真实复测的配对散点图及箱线图。}{候选模型在相同开发开局上的真实逐局得分。左侧连线对应同一开局，右侧展示箱线图与全部散点，红色菱形为均值。最佳模型在 28/30 局胜出，平均配对差为 137.87 分。}{fig:distribution}

\subsubsection{低分风险与配对差的进一步分析}
平均分提升之外，还需要考察策略是否减少了低分局。表\ref{tab:robustness}根据相同 30 个开发种子的逐局结果补充低分位数与低分局比例。分位数采用线性插值计算；以总分低于 50 分作为描述性阈值，该阈值用于比较两模型的低分尾部，不表示环境中的额外终止条件。

\begin{table}[!htbp]
 \centering\small
 \caption{相同 30 个开发种子上的波动与低分风险}\label{tab:robustness}
 \begin{tabularx}{.92\linewidth}{Y r r}
 \toprule
 指标 & 推荐最佳模型 & 续训最终模型\\
 \midrule
 得分标准差（总体形式） & 103.75 & 90.63\\
 第 10 百分位得分 & 111.90 & 19.90\\
 第 25 百分位得分 & 210.00 & 36.75\\
 低于 50 分的局数 & 1/30 & 9/30\\
 低于 50 分的比例 & 3.33\% & 30.00\%\\
 最低分--最高分 & 36--444 & 3--345\\
 时间截断局数 & 0 & 0\\
 完整通关局数（864 分） & 0 & 0\\
 \bottomrule
 \end{tabularx}
\end{table}

最佳模型的第 25 百分位得分为 210.00，续训最终模型为 36.75；低于 50 分的比例由 30.00\% 降至 3.33\%，相差 26.67 个百分点。这说明在所测开发开局上，提升同时体现在典型得分和低分尾部。最佳模型的绝对标准差仍达到 103.75，且没有一局完整通关，因此上述改进仍伴随局间波动，不能表述为稳定通关。

对每个开局计算配对差 $\Delta_i=R_i^{\mathrm{best}}-R_i^{\mathrm{final}}$，可进一步观察收益是否覆盖全部开局。图\ref{fig:robustness}显示，最佳模型在 28 局中得分较高，但在另外两局中分别少得 198 分和 9 分。均值差为 137.87 分，并不意味着每个开局都改善。对完整回合对反复重采样，得到均值差的百分位 Bootstrap 95\% 区间为 $[95.77,177.43]$ 分。该区间仅描述固定两模型在当前开发样本上的配对差变异，不覆盖独立重复训练的变异，也不替代最终留出测试。

\reportfigure{images/score_robustness.png}{低分风险与配对差}{实际逐局分数的经验累积分布及排序配对差。}{低分风险与逐种子收益。左侧为经验累积分布，虚线标记 50 分阈值；右侧为从小到大排序的最佳模型减最终模型得分，负值对应局部退步，红色虚线为平均配对差。}{fig:robustness}

\subsubsection{训练回退与恢复实验分析}
训练后期的回退首先反映出模型参数更新与策略性能之间并非单调关系。DQN 的学习目标包含目标网络估计，回放池又不断接收新策略产生的数据。即使当前小批量上的损失降低，也可能改变关键状态下的动作排序，进而影响之后整条轨迹。一次错误移动造成失球，还会减少后续继续获得奖励的机会。

回放分布变化是可能的解释之一。当回放容量有限时，较早阶段的有效经验会逐步被新经验替换；若新策略访问的状态变窄，网络可能减少对某些关键接球状态的训练。另一方面，自举估计误差、探索变化和不同开局下的轨迹差异也可能影响结果。现有汇总统计不能单独识别这些因素的贡献，因而回退原因仍属于需要通过日志和控制实验检验的解释。

三步回报恢复结束均分为 187.78，一步扩大回放结束均分为 216.94，均低于起点的 283.48。这说明在已执行设置下，两种恢复方案都没有形成超过起点的新最佳策略。三步目标引入更多真实奖励，并不保证在当前数据分布上优于一步目标；扩大回放也会同时改变经验新旧比例，不能简单等同于增强稳定性。

早停和最佳模型保存的作用是保留已有成果并减少继续退化阶段的计算支出。它们改变的是训练管理和模型选择方式，并没有使每一步更新都具备性能保证。当前选择依据应是独立评估结果，而不是模型文件的生成时间或累计训练步数。

\subsubsection{实时运行与行为分析}
实时游戏画面适合补充统计分数背后的行为解释。对同一局游戏，应按时间顺序记录小球下降、挡板移动、接触瞬间和反弹后的连续画面，同时保留模型、种子、分数及生命信息。仅有单帧画面难以判断运动方向，也不能说明一次接球是主动预测还是偶然重合。

\reportfigure{images/behavior_sequence.png}{模型行为连续帧}{同一模型、同一局中的成功接球与失球连续帧。}{同一局中的接球与失球对照，每排按相同时间间隔顺序采样。画面裁取下方接球区域，黄圈标记小球，蓝箭头标记挡板位移；子图保留实际得分与剩余生命，用于辅助判断运动变化。}{fig:behavior}

行为分析重点关注挡板与小球的相对运动。图\ref{fig:behavior}取自推荐最佳模型在同一局游戏中的两个片段。成功片段中，挡板在小球下降期间向接触位置移动，小球随后向上反弹。失球片段中，小球在右侧接近挡板高度时横向位置超出挡板覆盖范围，随后落到挡板下方，剩余生命减少。后续向左移动未能挽回这次失球。这些连续帧说明该模型仍会在局部轨迹上失误；单个成功片段不能证明其已学会精确落点预测，也不能据此推断所有失球的原因。

DQN 开发流程最终保留低学习率续训的最佳检查点。其与 PPO 的冻结模型测试统一见第\ref{sec:comparison}节，开发复测均分不与最终测试均分混合计算。

\FloatBarrier
